\chapter{Результаты экспериментов}
\label{cha:Results}

В данной главе представлены результаты вычислительных экспериментов.

\section{Тестовые данные}

Для тестирования использовались реальные данные от компаний — 436 тестов, каждый из которых представляет собой список коробок с различными размерами и количествами. Тестовые данные хранятся в формате CSV с полями:

\begin{verbatim}
SKU,Quantity,Length,Width,Height,Weight,Strength,Aisle,Caustic
\end{verbatim}

Размеры паллеты: $1200 \times 800$ мм (стандарт EURO).

\section{Условия эксперимента}

Эксперименты проводились на сервере со следующими характеристиками:
\begin{itemize}
    \item Процессор: 32 ядра
    \item Каждый тест запускался в отдельном потоке
    \item Язык реализации: C++20
    \item Компилятор: GCC с оптимизациями
\end{itemize}

\section{Сравнение алгоритмов}

\begin{table}[ht]
\centering
\caption{Сравнительные результаты алгоритмов}
\label{tab:benchmark}
\begin{tabular}{|l|c|c|c|c|c|}
\hline
\textbf{Алгоритм} & \textbf{Лимит} & \textbf{Среднее} & \textbf{Мин} & \textbf{Макс} & \textbf{Время} \\
\hline
GreedySolver & — & 73.56\% & 59.01\% & 85.21\% & 308 мс \\
LNSSolver & 1 сек & 77.18\% & 72.86\% & 86.39\% & 14 сек \\
LNSSolver & 5 сек & 78.23\% & 73.11\% & 88.15\% & 70 сек \\
LNSSolver & 30 сек & 78.93\% & 75.36\% & 88.65\% & 7 мин \\
LNSSolver & 300 сек & 79.71\% & 76.03\% & 90.39\% & 70 мин \\
\hline
\end{tabular}
\end{table}

\section{Анализ результатов}

\subsection{GreedySolver}

Жадный алгоритм демонстрирует высокую скорость работы — все 436 тестов обрабатываются за 308 миллисекунд. Однако качество укладки варьируется в широком диапазоне (59\%--85\%), что обусловлено сильной зависимостью от порядка коробок и их размеров.

Преимущества:
\begin{itemize}
    \item Очень высокая скорость
    \item Детерминированный результат
    \item Простота реализации
\end{itemize}

Недостатки:
\begin{itemize}
    \item Нестабильное качество на разных тестах
    \item Отсутствие оптимизации порядка укладки
\end{itemize}

\subsection{LNSSolver}

LNSSolver показывает значительно более стабильные результаты: минимальный коэффициент заполнения составляет 72.86\% даже при лимите в 1 секунду, что выше худшего результата GreedySolver.

Наблюдается логарифмическая зависимость качества от времени работы:
\begin{itemize}
    \item Первая секунда даёт прирост +3.6 п.п. относительно GreedySolver
    \item Увеличение времени в 5 раз (1с $\rightarrow$ 5с) даёт +1.05 п.п.
    \item Увеличение времени в 6 раз (5с $\rightarrow$ 30с) даёт +0.7 п.п.
    \item Увеличение времени в 10 раз (30с $\rightarrow$ 300с) даёт +0.78 п.п.
\end{itemize}

\subsection{Практические рекомендации}

\begin{enumerate}
    \item Для систем реального времени рекомендуется использовать LNSSolver с лимитом 1--5 секунд
    \item Для офлайн-планирования оптимален лимит 30--300 секунд
    \item GreedySolver может использоваться для быстрой оценки или как начальное приближение
\end{enumerate}

\section{Ограничения работы}

Текущая реализация имеет ряд ограничений, которые могут быть устранены в будущих версиях:

\subsection{Не учитываемые ограничения}

\begin{enumerate}
    \item \textbf{Хрупкость товаров}: Алгоритм не учитывает параметр Strength из входных данных. Хрупкие товары могут быть размещены под тяжёлыми.
    
    \item \textbf{Ограничение по весу на слой}: Не контролируется суммарный вес коробок, опирающихся на одну коробку.
    
    \item \textbf{Совместимость товаров}: Параметр Caustic не используется. Химически несовместимые товары могут быть размещены рядом.
    
    \item \textbf{Последовательность выгрузки}: Не оптимизируется порядок укладки для удобства выгрузки по секциям (Aisle).
    
    \item \textbf{Физическая устойчивость}: Метрики stability и center\_of\_mass\_z рассчитываются, но не используются в целевой функции оптимизации.
\end{enumerate}

\subsection{Технические ограничения}

\begin{enumerate}
    \item \textbf{Однопаллетная укладка}: Алгоритм работает только с одной паллетой. Многопаллетная укладка требует отдельной реализации.
    
    \item \textbf{Фиксированный размер паллеты}: Размер $1200 \times 800$ мм задан константой. Для других размеров требуется перекомпиляция.
    
    \item \textbf{Детерминированность}: GreedySolver детерминирован, LNSSolver использует случайный поиск и даёт разные результаты при повторных запусках.
\end{enumerate}

\subsection{Направления улучшения}

\begin{itemize}
    \item Добавить штраф за нарушение ограничений хрупкости в целевую функцию
    \item Реализовать проверку устойчивости с физическим симулятором
    \item Добавить многопаллетный режим с балансировкой нагрузки
    \item Интегрировать с системой управления складом (WMS API)
\end{itemize}

\section{Выводы}

\begin{enumerate}
    \item Разработанный алгоритм LNSSolver достигает коэффициента заполнения до 90\%, что соответствует современным промышленным решениям
    \item Жадный поиск в большой окрестности (LNS) с разнообразными операторами мутации эффективно исследует пространство решений, обеспечивая стабильное качество на всех тестах
    \item Структура HeightHandler обеспечивает эффективную работу с картой высот, что критически важно для производительности
    \item Визуализатор позволяет анализировать качество укладки и выявлять потенциальные улучшения алгоритма
\end{enumerate}
